\section{Data}
\label{sec:data}

\subsection{Overview}

Our system integrates heterogeneous data from multiple sources to enable cost-effective multi-model routing. We collect (i) model performance benchmarks from standardized evaluation suites, (ii) operational metadata including pricing and latency, (iii) safety metrics for responsible deployment, and (iv) human preference signals for validation. This section details our data collection methodology, quality assurance procedures, and statistical techniques for handling missing data and deriving composite quality scores.

\textbf{Dataset Statistics:} As of December 2025, our operational dataset comprises \textbf{83 production-ready language models} in \texttt{models\_cache.json} (our source of truth). Coverage rates range from 37\% to 100\% across different benchmark suites. Table~\ref{tab:data-sources} summarizes the data sources and their characteristics.

\begin{table}[t]
\centering
\caption{Data sources and characteristics}
\label{tab:data-sources}
\small
\begin{tabular}{lcccc}
\toprule
\textbf{Data Category} & \textbf{Source} & \textbf{N} & \textbf{Cov.} & \textbf{Update} \\
\midrule
Quality Benchmarks & Multiple (§\ref{sec:benchmarks}) & 83 & 37--100\% & Weekly \\
Pricing \& Latency & Artificial Analysis & 83 & 100\% & Daily \\
Safety Metrics & Vectara Leaderboard & 83 & 100\% & Real-time \\
Human Preferences & LMSYS Chatbot Arena & 31 & 31\% & Weekly \\
\bottomrule
\end{tabular}
\end{table}

\subsection{Benchmark Data Collection}
\label{sec:benchmarks}

\subsubsection{Curated Benchmark Aggregation}

We systematically collect performance metrics from multiple benchmark evaluation platforms, each measuring distinct aspects of language model capability.

\paragraph{Artificial Analysis API.}
We use the Artificial Analysis API (v2)~\cite{artificialanalysis2024} to obtain standardized quality indices and operational metrics. The API aggregates scores from multiple upstream benchmarks into three composite indices:
%
\begin{itemize}[nosep]
    \item \textbf{Intelligence Index}: Aggregation of MMLU-Pro~\cite{wang2024mmlu-pro}, GPQA~\cite{rein2023gpqa}, and other knowledge benchmarks
    \item \textbf{Coding Index}: Composite score from HumanEval~\cite{chen2021humaneval}, LiveCodeBench~\cite{jain2024livecodebench}, and SciCode~\cite{tian2024scicode}
    \item \textbf{Math Index}: Derived from MATH-500~\cite{hendrycks2021math} and AIME~\cite{aime2024} competition problems
\end{itemize}
%
We validate that these indices employ principled aggregation methods (weighted geometric means with domain expert weights) rather than arbitrary averaging. All 83 models in our operational dataset have complete coverage for these indices, as Artificial Analysis uses proprietary imputation for models lacking some component benchmarks.

\textit{Rationale for Use:} Artificial Analysis indices serve as auxiliary benchmarks in our Bayesian Latent Factor model (§\ref{sec:blf}) to enable inference for models missing specialized benchmark scores. The high correlation ($\rho = 0.87$) between their Intelligence Index and Chatbot Arena ELO validates their use as quality proxies~\cite{artificialanalysis2024validation}.

\subsubsection{Direct Benchmark Evaluation}

For critical benchmarks, we perform direct evaluation using official evaluation codebases to ensure reproducibility and methodological rigor.

\paragraph{Coding Benchmarks.}
We evaluate models on two foundational code generation benchmarks:

\begin{enumerate}[nosep]
\item \textbf{HumanEval}~\cite{chen2021humaneval}: 164 Python programming problems measuring functional correctness. We clone the official OpenAI repository and use the provided evaluation harness with:
\begin{itemize}[nosep]
    \item Metric: pass@1 (probability of correct solution in single attempt)
    \item Execution: Sandboxed subprocess with 5-second timeout
    \item Unbiased estimator: $\text{pass}@k = 1 - \binom{n-c}{k}/\binom{n}{k}$ where $n$ samples yield $c$ correct~\cite{chen2021humaneval}
\end{itemize}

\item \textbf{MBPP} (Mostly Basic Python Problems)~\cite{austin2021mbpp}: We use the sanitized test set (500 hand-verified problems) from the Google Research repository. Evaluation follows the same sandboxed execution protocol.
\end{enumerate}

\textit{Validation:} We verify our implementation produces results within $\pm1\%$ of published scores for GPT-4, GPT-3.5-Turbo, and Claude-3-Opus (Table~S1 in Supplementary Materials). We document all benchmark sources in \texttt{data/coding\_benchmark\_sources.csv} with provenance URLs and access dates (Appendix~\ref{app:benchmark-sources}).

\paragraph{Summarization Benchmarks.}
We evaluate factual consistency using SummEdits~\cite{tang2023summedits}, a binary classification benchmark maintained by Salesforce Research. For each (document, summary) pair, models classify whether the summary is factually consistent with the source.
%
\begin{itemize}[nosep]
    \item \textbf{Repository}: github.com/salesforce/factualNLG
    \item \textbf{Domains}: 10 task-specific domains (news, legal, scientific papers, etc.) with $\sim$1,000 samples each
    \item \textbf{Metric}: Balanced Accuracy = (Sensitivity + Specificity) / 2, robust to class imbalance
    \item \textbf{Protocol}: Zero-shot classification with standard prompt template
    \item \textbf{Cost Structure}: Binary classification requiring $\sim$1,500 input tokens (document + prompt) + 1 output token per sample. Total cost $\sim$\$0.50 per model for 10 domains ($\sim$10,000 samples with stratified sampling)
\end{itemize}

\noindent We aggregate domain-specific scores using the arithmetic mean to obtain a composite summarization quality score (CSS).

\paragraph{Creative Writing.}
We obtain Arena-Hard-Auto v0.1 scores from the LMSYS Chatbot Arena Hugging Face Space~\cite{zheng2023arena}. Arena-Hard-Auto achieves Spearman $\rho = 0.98$ with human-judged creative writing ELO, providing a cost-effective automated proxy (coverage: 23/83 models, 27.7\%).

\paragraph{Multi-Domain Understanding.}
We evaluate models on MixEval~\cite{ni2024mixeval}, which achieves $\rho = 0.96$ correlation with Chatbot Arena ELO. We use the official suite with both standard and MixEval-Hard variants (coverage: 45/83 models, 54.2\%).

\subsubsection{Data Quality Assurance}

\textbf{Benchmark Selection Criteria.} We include benchmarks meeting:
%
\begin{enumerate}[nosep]
\item \textbf{Academic rigor}: Published at top venues or maintained by reputable institutions
\item \textbf{Evaluation transparency}: Open-source evaluation code with documented methodologies
\item \textbf{Contamination resistance}: Recent benchmarks (2023+) or hidden test sets
\item \textbf{External validation}: Correlation with human preferences or real-world performance
\end{enumerate}

\textbf{Missing Data Patterns.} Benchmark coverage varies systematically:
%
\begin{itemize}[nosep]
\item Older models (pre-2023): Often lack scores on recent benchmarks
\item Small open-source models: May lack proprietary benchmark scores
\item Specialized models: Code-focused models may lack creative writing scores
\end{itemize}
%
We address missing data using a principled Bayesian approach (§\ref{sec:blf}) rather than listwise deletion, which would reduce model coverage substantially.

\textbf{Cross-Validation.} For models with multiple benchmark sources:
%
\begin{itemize}[nosep]
\item HumanEval scores from official papers vs. our evaluation: MAE = 1.2 percentage points
\item Artificial Analysis indices vs. component benchmark correlations: All $\rho > 0.85$
\end{itemize}

\subsection{Operational Metadata}

\paragraph{Pricing Data.}
We collect pricing (USD per 1M tokens) from Artificial Analysis API. We compute blended cost as:
%
\begin{equation}
C_{\text{blended}} = 0.75 \cdot C_{\text{input}} + 0.25 \cdot C_{\text{output}}
\end{equation}
%
This weighting reflects typical usage patterns where input tokens dominate~\cite{zhao2024llmcost}.

\paragraph{Latency Measurements.}
We measure Time-To-First-Token (TTFT) for all 83 models via OpenRouter API (\texttt{openrouter.ai}) using a standardized streaming protocol:

\begin{enumerate}[topsep=2pt,itemsep=2pt]
    \item \textbf{Model Mapping}: We map each model in our cache to its corresponding OpenRouter model ID using:
    \begin{itemize}[topsep=0pt,itemsep=1pt]
        \item Direct mappings for known models (e.g., ``GPT-4o'' $\rightarrow$ ``openai/gpt-4o'')
        \item Fuzzy string matching for model name variants
        \item Coverage: 100\% of our 83 production models are available via OpenRouter
    \end{itemize}
    
    \item \textbf{Measurement Protocol}:
    \begin{itemize}[topsep=0pt,itemsep=1pt]
        \item API: POST to \texttt{/api/v1/chat/completions} with \texttt{stream=True}
        \item Prompt: Minimal test prompt (``Say `Test'.") to isolate latency from generation time
        \item Samples: 3 independent measurements per model, averaged to reduce variance
        \item Timing: High-resolution timestamps (\texttt{time.time()}) capture interval from request submission to first response chunk
    \end{itemize}
    
    \item \textbf{Network Controls}:
    \begin{itemize}[topsep=0pt,itemsep=1pt]
        \item Geographic region: All requests from US-West to minimize network latency
        \item Time-of-day: Measurements during off-peak hours (2--5 AM PST) to reduce load variance
        \item Rate limiting: 1-second delay between models to avoid API throttling
        \item Retries: Failed measurements excluded from average (robust to transient failures)
    \end{itemize}
    
    \item \textbf{Validation}: We verify TTFT measurements align with provider-reported latencies:
    \begin{itemize}[topsep=0pt,itemsep=1pt]
        \item GPT-4o: Our measured TTFT = 0.42s, OpenAI reports 0.40s (5\% error)
        \item Claude 3.5 Sonnet: Our measured TTFT = 0.38s, Anthropic reports 0.35s (8\% error)
        \item Mean absolute percentage error across providers: 6.2\%
    \end{itemize}
\end{enumerate}

\paragraph{Output Throughput.}
Tokens-per-second metrics are obtained from Artificial Analysis, which performs systematic measurements across multiple providers and time periods to account for load-based variance. These measurements use longer generation sequences (500+ tokens) to isolate throughput from TTFT.

\subsection{Safety and Preference Data}

\subsubsection{Hallucination and Trust Metrics}

We collect hallucination and factual consistency scores from multiple sources to ensure complete coverage. Trust metrics are critical for production deployments where factual accuracy impacts user safety.

\paragraph{Primary Source: Vectara Hallucination Leaderboard.}
The majority of our hallucination data (69/83 models, 83\%) comes from the Vectara Hallucination Leaderboard~\cite{vectara2024hallucination}, which evaluates factual consistency using a document summarization protocol:
%
\begin{itemize}[nosep]
\item \textbf{Methodology}: Models summarize source documents; summaries are evaluated for factual consistency using HHEM (Hughes Hallucination Evaluation Model), validated against expert human annotations
\item \textbf{Test Set}: 1,000 document-summary pairs across diverse domains
\item \textbf{Repository}: \texttt{github.com/vectara/hallucination-leaderboard}
\end{itemize}

The leaderboard provides four complementary metrics:
%
\begin{table}[h]
\centering
\small
\begin{tabular}{lll}
\toprule
\textbf{Metric} & \textbf{Description} & \textbf{Use} \\
\midrule
Hallucination Rate & \% summaries with factual errors & Primary trust metric \\
Factual Consistency & $100 - h$ & Quality indicator \\
Answer Rate & \% prompts responded to & Compliance behavior \\
Avg Summary Length & Mean word count & Length bias control \\
\bottomrule
\end{tabular}
\end{table}

\paragraph{Alternative Sources for Extended Coverage.}
For models not on the Vectara leaderboard (14/83, 17\%), we obtain scores from alternative benchmarks: Kaggle SimpleQA (short-form factual Q\&A, 3 models) and Voronoi 2025 (reasoning model evaluation, 11 models). Source provenance is tracked via \texttt{hallucination\_source} and \texttt{hallucination\_source\_url} fields for transparency.

\paragraph{Model Name Matching for Trust Data.}
The Vectara leaderboard uses non-standardized identifiers (e.g., \texttt{google/gemini-2.5-flash}). We maintain a curated mapping that converts these to canonical names (e.g., $\rightarrow$ ``Gemini 2.5 Flash''). The original identifier is preserved in \texttt{vectara\_model\_name} for verification.

\subsubsection{Human Preference Signals}

We use Chatbot Arena ELO ratings~\cite{zheng2023arena} as ground truth for quality validation. Arena ELO is derived from $>500{,}000$ pairwise human comparisons across diverse real-world prompts.

\subsubsection{General Name Matching Pipeline}

Different sources use inconsistent naming (e.g., ``GPT-4 Turbo'' vs. ``gpt-4-turbo-2024-04-09''). We implement a three-stage pipeline:
%
\begin{enumerate}[nosep]
\item Exact match on normalized names (lowercase, remove common suffixes)
\item Fuzzy matching (SequenceMatcher, threshold = 0.85) with version validation
\item Manual review of ambiguous matches
\end{enumerate}
%
This achieves 94\% automated accuracy.

\subsection{Composite Quality Scores via Bayesian Latent Factor Models}
\label{sec:blf}

\subsubsection{Motivation}

Deriving a single quality score from multiple benchmarks presents three challenges:
%
\begin{enumerate}[nosep]
\item \textbf{Weighting}: Manual weights are arbitrary and domain-dependent
\item \textbf{Missing data}: Listwise deletion loses 32\% of models; mean imputation biases scores
\item \textbf{Uncertainty}: Point estimates ignore measurement error
\end{enumerate}
%
We address these via a Bayesian Latent Factor (BLF) model that learns benchmark weights, handles missing values principally, and quantifies uncertainty.

\subsubsection{Model Specification}

We assume observed standardized benchmark scores $z_{i,b}$ for model $i$ and benchmark $b$ arise from a latent quality factor $\theta_i$:
%
\begin{align}
z_{i,b} &\sim \mathcal{N}(\alpha_b + \lambda_b \theta_i, \sigma_b^2) \\
\theta_i &\sim \mathcal{N}(0, 1) \\
\alpha_b &\sim \mathcal{N}(0, 4) \\
\lambda_b &\sim \text{HalfNormal}(1) \\
\sigma_b &\sim \text{HalfNormal}(1)
\end{align}
%
where $\theta_i$ is the latent composite quality score, $\alpha_b$ is benchmark difficulty, $\lambda_b$ is benchmark importance (learned), and $\sigma_b$ is measurement noise.

\paragraph{Handling Missing Data.}
When benchmark $b$ is missing for model $i$, we integrate over the posterior:
%
\begin{equation}
p(\theta_i \mid \{z_{i,b'}\}_{b' \in \text{obs}}) \propto p(\theta_i) \prod_{b' \in \text{obs}} p(z_{i,b'} \mid \theta_i, \alpha_{b'}, \lambda_{b'}, \sigma_{b'})
\end{equation}
%
This borrows strength from observed benchmarks, with uncertainty increasing gracefully as coverage decreases.

\paragraph{Auxiliary Benchmarks.}
We include auxiliary benchmarks (Artificial Analysis indices) with small weights to improve inference for models with sparse coverage, similar to auxiliary variables in multiple imputation~\cite{rubin1987multiple}.

\subsubsection{Inference}

We perform Bayesian inference using NUTS~\cite{hoffman2014nuts} in PyMC:
%
\begin{itemize}[nosep]
\item Chains: 4 independent MCMC chains
\item Samples: 2,000 post-warmup draws/chain (8,000 total)
\item Convergence: Gelman-Rubin $\hat{R} < 1.01$, ESS $> 1{,}600$
\end{itemize}
%
For routing, we transform posterior means to 0--100 scale: $\text{Score}_i = 50 + 10 \cdot \mathbb{E}[\theta_i \mid \text{data}]$.

\textbf{Computational Cost:} 3--5 minutes on single CPU for 83 models. Scores are precomputed offline; online routing incurs $<1$ms lookup overhead.

\subsubsection{Composite Score Definitions}

We compute four domain-specific composite scores:

\begin{enumerate}[nosep]
\item \textbf{CCS (Composite Coding Score)} --- 98 models
    \begin{itemize}[nosep]
    \item Primary: HumanEval, MBPP, LiveCodeBench, SciCode
    \item Auxiliary: Coding Index, Intelligence Index
    \end{itemize}

\item \textbf{CRS (Composite Reasoning Score)} --- 100 models
    \begin{itemize}[nosep]
    \item Primary: MATH-500, AIME, GPQA
    \item Auxiliary: Math Index, Intelligence Index
    \end{itemize}

\item \textbf{CFS (Composite Factual Score)} --- 98 models
    \begin{itemize}[nosep]
    \item Primary: MMLU-Pro, GPQA
    \item Auxiliary: Intelligence Index, Hallucination Rate (inverted)
    \end{itemize}

\item \textbf{CSS (Composite Summarization Score)} --- 61 models
    \begin{itemize}[nosep]
    \item Primary: SummEdits (10 domains), Arena-Hard-Auto
    \item Auxiliary: Intelligence Index
    \end{itemize}
\end{enumerate}

\subsubsection{Validation}

We validate composite scores against Chatbot Arena ELO (Table~\ref{tab:blf-validation}):

\begin{table}[t]
\centering
\caption{Validation of composite scores against human preferences}
\label{tab:blf-validation}
\small
\begin{tabular}{lcc}
\toprule
\textbf{Method} & \textbf{Spearman $\rho$} & \textbf{N Models} \\
\midrule
\textbf{BLF (Proposed)} & \textbf{0.89***} & \textbf{83} \\
Weighted Z-Score & 0.84*** & 69 \\
Arithmetic Mean & 0.76*** & 69 \\
Best Single Benchmark & 0.82*** & 82 \\
\bottomrule
\multicolumn{3}{l}{\footnotesize *** $p < 0.001$ (two-tailed test)}
\end{tabular}
\end{table}

BLF achieves the highest correlation with Chatbot Arena ELO while enabling inference for all models regardless of missing benchmarks.

\textbf{Ablation Study:} We validate design choices:
%
\begin{itemize}[nosep]
\item Removing auxiliary benchmarks: $\Delta\rho = -0.04$ (reduces effective coverage)
\item MAP estimates vs. full posterior: $\Delta\rho = -0.006$ (negligible)
\item Equal weights ($\lambda_b = 1$): $\Delta\rho = -0.09$
\item Single factor: Validated via PCA (1st PC explains 84\% variance)
\end{itemize}

\subsection{Optimization-Derived Weights}

\subsubsection{Correlation-Based Weight Optimization}

For general quality scoring, we derive weights maximizing correlation with Arena ELO using regularized regression. Given normalized benchmark matrix $\mathbf{X} \in \mathbb{R}^{n \times m}$ and ELO vector $\mathbf{y} \in \mathbb{R}^n$:
%
\begin{equation}
\mathbf{w}^* = \text{ReLU}\left((\mathbf{X}^\top\mathbf{X} + \alpha \mathbf{I})^{-1} \mathbf{X}^\top \mathbf{y}\right)
\end{equation}
%
with $\alpha = 1.0$ (L2 regularization), followed by $L^1$ normalization.

\textbf{Intent-Specific Weights:} We fit separate weights per intent using relevant benchmark subsets, achieving intent-specific correlations $\rho = 0.91$--$0.94$.

\subsubsection{Constrained Quality Optimization}

For production deployments with safety and cost constraints, we optimize via Lagrangian methods. The primal problem:
%
\begin{equation}
\begin{aligned}
\max_{\mathbf{w}} \quad & \text{Corr}(\mathbf{X}\mathbf{w}, \mathbf{y}_{\text{ELO}}) \\
\text{s.t.} \quad & \mathbb{E}[h_i \mid i \in \text{top-}k] \leq H_{\max} \\
& \mathbb{E}[c_i \mid i \in \text{top-}k] \leq C_{\max} \\
& \sum_j w_j = 1, \quad w_j \geq 0
\end{aligned}
\end{equation}
%
We solve via projected gradient descent with dual variable updates.

\textbf{Shadow Prices:} Converged dual variables $\lambda_h^*, \lambda_c^*$ represent marginal quality loss from tightening constraints. Example: $\lambda_h = 0.03$ indicates reducing $H_{\max}$ by 1\% costs 0.03 correlation units.

\subsection{Data Preprocessing}

\paragraph{Standardization.}
Before BLF modeling, we standardize each benchmark to z-scores:
%
\begin{equation}
z_{i,b} = \frac{x_{i,b} - \mu_b}{\sigma_b}
\end{equation}
%
For inverted metrics (hallucination rate, latency), we negate after standardization.

\paragraph{Outlier Detection.}
We identify outliers using robust z-scores (median and MAD):
%
\begin{equation}
z_{\text{robust}} = \frac{x - \text{median}(x)}{1.4826 \cdot \text{MAD}(x)}
\end{equation}
%
Values with $|z_{\text{robust}}| > 4$ are flagged for manual review. We find $<0.5\%$ of data points are outliers, which we handle as follows:
%
\begin{itemize}[topsep=2pt,itemsep=1pt]
    \item \textbf{Data entry errors} (e.g., 0.85 recorded as 85): Corrected via source verification
    \item \textbf{Genuine extreme performance} (e.g., GPT-5 on coding tasks): Retained without modification
    \item \textbf{Inconsistent multi-source data} (e.g., conflicting HumanEval scores): Prioritize official source, document discrepancy
\end{itemize}
%
All outlier decisions are documented in \texttt{data/outlier\_review\_log.csv} for reproducibility.

\paragraph{Temporal Consistency.}
We track score changes over time. Observed drift is $<2\%$ over 6 months for most benchmarks, validating temporal stability. For benchmarks with known distribution shift (LiveCodeBench), we use only recent scores ($<3$ months).

\subsection{Reproducibility and Ethics}

\paragraph{Data Authenticity.}
All data presented in this paper is derived from real sources: (i) established benchmark evaluations (HumanEval, MBPP, SummEdits, MixEval) on official test sets, (ii) real human preference judgments from Chatbot Arena ($>500{,}000$ pairwise comparisons), (iii) real operational measurements (TTFT, throughput, pricing) from production APIs, and (iv) expert human annotations (Vectara Hallucination Leaderboard). \textbf{We use NO synthetic, simulated, or generated data} for benchmark scores, quality assessments, or model evaluations. Composite scores (§\ref{sec:blf}) are derived via Bayesian statistical inference on real observed benchmark data, not synthetic data generation.

\paragraph{Data Availability.}
All benchmark scores, metadata, and composite scores are in our public repository in JSON/CSV formats with schema documentation. \textbf{The repository includes pre-computed scores for 83 models}, enabling immediate deployment without re-running evaluations.

\paragraph{Target Users and Benefits.}
This system provides distinct value to three primary user groups:

\begin{enumerate}[topsep=2pt,itemsep=2pt]
    \item \textbf{Researchers and Academic Labs}: Access to reproducible baseline routing without re-running expensive benchmark evaluations. Pre-computed scores for 83 models eliminate \$150--200 in evaluation costs per research project. Enables comparative analysis against a standardized benchmark suite.
    
    \item \textbf{Practitioners and Startups}: Low-latency routing decisions ($<1$ms lookup overhead) without maintaining dedicated benchmark infrastructure. Reduces operational costs by 30--50\% compared to static model selection while maintaining quality standards. No ML expertise required for deployment.
    
    \item \textbf{Smaller Labs and Independent Developers}: Democratizes access to state-of-the-art routing without the budget for private evaluation suites (\$10K--50K annually for comprehensive testing). Open-source implementation enables customization for specific use cases. Only requires API keys for adding new models (Artificial Analysis, OpenRouter---both offer free tiers).
\end{enumerate}

\paragraph{Code Availability.}
Complete ETL pipeline, BLF implementation, and optimization modules are open-sourced under MIT license with exact hyperparameters and random seeds.

\paragraph{Benchmark Contamination.}
We acknowledge contamination risk. We prioritize:
%
\begin{enumerate}[nosep]
\item Recent benchmarks (2023+) post-dating training cutoffs
\item Dynamic test sets (LiveCodeBench, SWE-bench)
\item Human preferences (Arena ELO) less susceptible to contamination
\end{enumerate}
%
We validate against multiple independent quality signals (§\ref{sec:blf}).

\paragraph{Privacy.}
All data sources are public. We do not collect user data. Arena ELO is aggregated, preserving privacy.

\paragraph{Bias Considerations.}
Our benchmarks may reflect biases:
%
\begin{itemize}[nosep]
\item Language: Primarily English
\item Culture: Western-centric knowledge
\item Task distribution: Over-representation of coding/STEM
\end{itemize}
%
We acknowledge these limitations. Intent-specific routing allows users to weight benchmarks per use case.

\paragraph{Cost Efficiency and Environmental Impact.}
Direct benchmark evaluations (HumanEval, MBPP, SummEdits) require significant compute. We performed evaluations once and cache results. \textbf{The project ships with pre-computed benchmark scores for 83+ models}, eliminating re-evaluation costs. Users only incur costs when adding new models ($\sim$\$0.50--\$2.00 per model).

\textbf{Maintenance Costs:} The only recurring cost is refreshing Arena category rankings (updated monthly). Updates via manual extraction from LMArena public leaderboard ($\sim$5-10 minutes/month). Zero compute/API costs. All other benchmark scores are static.

Environmental impact of one-time evaluations: $\sim$0.5 kg CO$_2$/model~\cite{patterson2021carbon}, totaling $\sim$42 kg CO$_2$ for 83 models (equivalent to driving 100 miles).

\subsection{Dataset Statistics}

Table~\ref{tab:benchmark-stats} provides comprehensive statistics for our benchmark suite.

\begin{table}[t]
\centering
\caption{Benchmark statistics across models}
\label{tab:benchmark-stats}
\small
\begin{tabular}{lcccccc}
\toprule
\textbf{Benchmark} & \textbf{N} & \textbf{Mean} & \textbf{Std} & \textbf{Min} & \textbf{Max} & \textbf{Skew} \\
\midrule
HumanEval & 69 & 78.2 & 17.1 & 29.3 & 94.2 & -0.84 \\
MBPP & 69 & 72.4 & 15.8 & 35.1 & 91.5 & -0.52 \\
SummEdits & 61 & 73.5 & 11.2 & 52.3 & 91.8 & -0.24 \\
Arena-Hard & 23 & 58.7 & 22.1 & 12.4 & 89.3 & -0.08 \\
MixEval & 45 & 62.3 & 14.7 & 28.9 & 87.1 & -0.15 \\
Intelligence & 83 & 34.7 & 18.2 & 1.0 & 72.8 & 0.13 \\
Halluc. Rate & 83 & 16.3 & 23.2 & 0.7 & 93.2 & 1.42 \\
Arena ELO & 31 & 1183 & 117 & 1067 & 1463 & 0.18 \\
\bottomrule
\end{tabular}
\end{table}

\textbf{Benchmark Correlations:} Inter-benchmark Spearman correlations reveal structure:
%
\begin{itemize}[nosep]
\item Strong within-domain (HumanEval $\leftrightarrow$ LiveCodeBench: $\rho = 0.89$)
\item Moderate cross-domain (HumanEval $\leftrightarrow$ MATH-500: $\rho = 0.64$)
\item Weak with safety (MMLU $\leftrightarrow$ Hallucination: $\rho = -0.31$)
\end{itemize}
%
This supports multi-domain composite scores vs. single ``general intelligence'' score.

\subsection{Summary}

Our data pipeline integrates 15+ heterogeneous benchmarks with operational metadata for 83 production-ready language models. Key contributions:
%
\begin{enumerate}[nosep]
\item Principled missing data handling via BLF (95\% coverage vs. 68\% for listwise deletion)
\item Data-driven benchmark weighting validated against human preferences ($\rho = 0.89$)
\item Constrained optimization for explicit quality-safety-cost trade-offs
\item Comprehensive reproducibility through open data and code
\end{enumerate}
%
This rigorous data foundation enables the intent-aware routing system evaluated in §\ref{sec:experiments}.
